\section{Machine Learning and Perception}

Likely examiners: Prof.\ Klambauer or Prof.\ Hochreiter. The same core comes back almost every
time: empirical risk minimization, gradient descent, backpropagation and the delta error. Then
short questions on activations, regularization, optimizers and one architecture. Notation of the
JKU script: input $\mathbf{x}_n$, label $y_n$, model $g(\mathbf{x};\mathbf{w})$, loss $L$,
pre-activation $s$, activation $a = f(s)$, layer $[l]$.

\subsection{Empirical risk minimization and generalization}

\begin{qabox}{Write down a supervised data set and the training framework.}
\origin{asked}{Klambauer (several students), Berger-Weber committee}
\core{We have $N$ pairs of input and label, and learning means minimizing the average loss on them.}
\explain The pairs are drawn independently from an unknown distribution. A model maps an input to a
prediction, and a loss says how wrong one prediction is. The average loss on the training data is
the empirical risk.
\derive
\begin{align*}
Z &= \{(\mathbf{x}_n,y_n)\}_{n=1}^{N} \why{training data, i.i.d. from $p(\mathbf{x},y)$} \\
\Remp(\mathbf{w}) &= \frac1N\sum_{n=1}^{N} L\bigl(y_n, g(\mathbf{x}_n;\mathbf{w})\bigr) \why{average loss} \\
\hat{\mathbf{w}} &= \arg\min_{\mathbf{w}} \Remp(\mathbf{w}) \why{learning}
\end{align*}
\summary{Data, model, loss, average, minimize.}
\src{DLNN script WS24, Ch.~3 and 5.1.}
\end{qabox}

\begin{qabox}{What is the generalization error? Why can we not minimize it directly?}
\origin{asked}{Klambauer (Fürnkranz committee), Schedl committee}
\core{It is the expected loss on new data, and we cannot compute it because the true distribution is unknown.}
\explain So we minimize the empirical risk instead and estimate the generalization error on data we
did not train on: a test set or cross-validation. The training loss is not a fair estimate, because
the weights were chosen on that data.
\derive
\begin{align*}
\Rgen(\mathbf{w}) &= \mathbb{E}_{(\mathbf{x},y)\sim p}\bigl[L(y,g(\mathbf{x};\mathbf{w}))\bigr] \why{expected loss} \\
&= \int L\bigl(y,g(\mathbf{x};\mathbf{w})\bigr)\,p(\mathbf{x},y)\,\mathrm{d}\mathbf{x}\,\mathrm{d}y \why{needs $p$, which is unknown} \\
&\approx \frac{1}{M}\sum_{m=1}^{M} L\bigl(y_m,g(\mathbf{x}_m;\mathbf{w})\bigr) \why{test set, law of large numbers}
\end{align*}
\summary{Risk = loss on unseen data. Estimate it with a test set or cross-validation.}
\src{DLNN script WS24, Ch.~5.2--5.3.}
\end{qabox}

\begin{qabox}{What is a loss function? Give examples.}
\origin{asked}{Klambauer, March 2026 (``not all losses are differentiable, 0--1 loss'')}
\core{A loss compares a label with a prediction and is zero when the prediction is right.}
\explain It is never negative. We choose it to fit the task. Not every loss can be trained with
gradient descent: the 0--1 loss has no useful gradient, so we use cross-entropy instead.
\explain Examples: squared loss $\tfrac12(y-\hat y)^2$ for regression; binary cross-entropy
$-y\log\hat y - (1-y)\log(1-\hat y)$; categorical cross-entropy $-\sum_k y_k\log\hat y_k$; hinge
loss $\max(0, 1 - y\,g(\mathbf{x}))$ for the SVM.
\summary{Non-negative, zero when right. 0--1 loss $\to$ use cross-entropy.}
\src{DLNN script WS24, Sec.~4.1--4.4.}
\end{qabox}

\begin{qabox}{Show ERM on linear regression.}
\origin{asked}{Klambauer (Küng committee), Schedl committee}
\core{With a linear model and squared loss the risk is convex, so the minimum has a closed form.}
\explain Set the gradient to zero and solve. Gradient descent reaches the same point step by step.
For one sample the gradient is error times input, and the same pattern comes back in backprop.
\derive
\begin{align*}
\Remp(\mathbf{w}) &= \tfrac12\lVert \mathbf{y} - X\mathbf{w}\rVert^2 \why{squared loss, all samples} \\
\nabla_{\mathbf{w}}\Remp &= -X^{\mathsf{T}}(\mathbf{y} - X\mathbf{w}) \why{chain rule} \\
0 &= X^{\mathsf{T}}\mathbf{y} - X^{\mathsf{T}}X\,\hat{\mathbf{w}} \why{gradient zero at the minimum} \\
\hat{\mathbf{w}} &= (X^{\mathsf{T}}X)^{-1}X^{\mathsf{T}}\mathbf{y} \why{normal equations, pseudo-inverse}
\end{align*}
\summary{Convex $\to$ closed form. Gradient = error $\times$ input.}
\src{DLNN script WS24, Sec.~4.1.3--4.1.4.}
\end{qabox}

\begin{qabox}{Overfitting, underfitting, bias and variance.}
\origin{predicted}{DLNN script Sec.~5.4; bias and MSE were asked only in signal processing (Huemer)}
\core{Too simple a model has high bias, too complex a model has high variance.}
\explain An underfitting model cannot follow the data. An overfitting model follows the noise, so the
training error is low and the test error is high. Very large networks show double descent: after
the point where they fit the training data exactly, the test error falls again.
\derive
\begin{align*}
y &= f(\mathbf{x}) + \varepsilon,\quad \mathbb{E}[\varepsilon] = 0,\ \mathrm{Var}[\varepsilon] = \sigma^2 \why{data model} \\
\mathbb{E}\bigl[(y-\hat g)^2\bigr] &= \mathbb{E}\bigl[(f + \varepsilon - \bar g + \bar g - \hat g)^2\bigr] \why{add and subtract $\bar g = \mathbb{E}[\hat g]$} \\
&= \sigma^2 + (f - \bar g)^2 + \mathbb{E}\bigl[(\hat g - \bar g)^2\bigr] \why{cross terms have zero mean} \\
&= \text{noise} + \text{bias}^2 + \text{variance}
\end{align*}
\summary{Simple $\to$ bias. Complex $\to$ variance. Large nets: double descent.}
\src{DLNN script WS24, Sec.~5.4.}
\end{qabox}

\subsection{From one neuron to a deep network}

\begin{qabox}{Write the formula of a neuron.}
\origin{asked}{Schlüter}
\core{A neuron takes a weighted sum of its inputs plus a bias, and passes it through a non-linear function.}
\explain The weighted sum is the pre-activation, and the result is the activation. A layer does
this for many neurons at once. The forward pass goes layer by layer from the input to the output.
\derive
\begin{align*}
s_i &= \sum_j w_{ij}\,a_j + b_i \why{pre-activation} \\
a_i &= f(s_i) \why{activation} \\
\mathbf{a}^{[l]} &= f\bigl(W^{[l]}\mathbf{a}^{[l-1]} + \mathbf{b}^{[l]}\bigr),\quad \mathbf{a}^{[0]} = \mathbf{x} \why{one layer; output $\hat{\mathbf{y}} = \mathbf{a}^{[L]}$}
\end{align*}
\summary{Weighted sum, plus bias, through a non-linearity.}
\src{DLNN script WS24, Eqs.~(3.2)--(3.3), (4.59)--(4.62).}
\end{qabox}

\begin{qabox}{A deep network with and without activation functions?}
\origin{asked}{Schlüter}
\core{Without non-linear activations the whole network collapses to one linear model.}
\explain A product of linear maps is again linear. So depth adds nothing, and the network cannot
even solve XOR. The non-linearity is what makes depth useful.
\derive
\begin{align*}
\mathbf{a}^{[2]} &= W^{[2]}\bigl(W^{[1]}\mathbf{x} + \mathbf{b}^{[1]}\bigr) + \mathbf{b}^{[2]} \why{two layers, no activation} \\
&= \bigl(W^{[2]}W^{[1]}\bigr)\mathbf{x} + \bigl(W^{[2]}\mathbf{b}^{[1]} + \mathbf{b}^{[2]}\bigr) \why{multiply out} \\
&= W\mathbf{x} + \mathbf{b} \why{one linear layer}
\end{align*}
\summary{No activation $\to$ one linear model. No XOR.}
\src{DLNN script WS24, Sec.~4.5.}
\end{qabox}

\begin{qabox}{Universal function approximation theorem.}
\origin{asked}{Lehner}
\core{One hidden layer with enough units can approximate any continuous function on a compact set.}
\explain The activation must be non-polynomial, for example a sigmoid. The theorem says the network
can represent the function. It does not say how many units we need, and it does not say that
gradient descent will find the weights.
\summary{One hidden layer, any continuous function. Existence, not learnability.}
\src{DLNN I+II script SS25, ``Deep Learning theory''.}
\end{qabox}

\subsection{Backpropagation and the delta error}

This is the most asked topic. Practise writing it until it takes two minutes.

\begin{qabox}{Explain backpropagation and define the delta error.}
\origin{asked}{Hochreiter, Klambauer, Kofler (several students)}
\core{Backprop gets all gradients with one forward and one backward pass, and the delta is the
derivative of the loss with respect to a pre-activation.}
\explain First the forward pass stores all $s$ and $a$. Then the deltas are computed from the
output back to the input. The gradient of a weight is the delta of the unit above times the
activation of the unit below. The cost is linear in the number of weights.
\derive
\begin{align*}
\delta_i &:= \frac{\partial L}{\partial s_i} \why{definition} \\
\frac{\partial L}{\partial w_{ij}} &= \frac{\partial L}{\partial s_i}\,\frac{\partial s_i}{\partial w_{ij}} \why{chain rule: $w_{ij}$ acts only through $s_i$} \\
&= \delta_i\,a_j \why{$s_i = \sum_j w_{ij}a_j + b_i$} \\
w_{ij} &\leftarrow w_{ij} - \eta\,\delta_i\,a_j \why{gradient step}
\end{align*}
\summary{$\delta$ = dLoss/d(pre-activation). Gradient = $\delta$ above $\times$ activation below.}
\src{DLNN script WS24, Sec.~4.6.2, Eqs.~(4.63)--(4.68).}
\end{qabox}

\begin{qabox}{Write the recursive formula for the deltas.}
\origin{asked}{Klambauer, Kofler (``delta error of the $l$-th layer'')}
\core{The delta of a hidden unit is the weighted sum of the deltas above, times the slope of its own activation.}
\explain A hidden unit $j$ affects the loss only through the units $i$ of the next layer. So we sum
over those paths. In matrix form we multiply by the weight matrix and then by the derivatives of
the activations.
\derive
\begin{align*}
\delta_j &= \frac{\partial L}{\partial s_j} = \sum_i \frac{\partial L}{\partial s_i}\,\frac{\partial s_i}{\partial a_j}\,\frac{\partial a_j}{\partial s_j} \why{chain rule over all units $i$ above} \\
&= \sum_i \delta_i\,w_{ij}\,f'(s_j) \why{$\partial s_i/\partial a_j = w_{ij}$, $\partial a_j/\partial s_j = f'(s_j)$} \\
\bm{\delta}^{[l-1]} &= \mathrm{diag}\bigl(f'(\mathbf{s}^{[l-1]})\bigr)\,W^{[l]\mathsf{T}}\,\bm{\delta}^{[l]} \why{layer form} \\
\frac{\partial L}{\partial W^{[l]}} &= \bm{\delta}^{[l]}\,\mathbf{a}^{[l-1]\mathsf{T}} \why{outer product}
\end{align*}
\summary{Deltas flow down: multiply by $W^{\mathsf T}$, then by $f'$.}
\src{DLNN script WS24, Eqs.~(4.71), (4.77)--(4.84).}
\end{qabox}

\begin{qabox}{What is the delta at the output with softmax and cross-entropy?}
\origin{asked}{Hochreiter (softmax + cross-entropy)}
\core{It is simply prediction minus label.}
\explain The same clean result holds for a linear output with squared loss and for a sigmoid with
binary cross-entropy. That is why these pairs are used together.
\derive
\begin{align*}
a_k &= \frac{e^{s_k}}{\sum_m e^{s_m}},\quad L = -\sum_k y_k\log a_k \why{softmax and cross-entropy} \\
\frac{\partial a_k}{\partial s_i} &= a_k\,(\mathbb{1}[k=i] - a_i) \why{derivative of the softmax} \\
\frac{\partial L}{\partial s_i} &= -\sum_k \frac{y_k}{a_k}\,a_k\,(\mathbb{1}[k=i] - a_i) \why{chain rule} \\
&= -y_i + a_i\sum_k y_k = a_i - y_i \why{one-hot label: $\sum_k y_k = 1$}
\end{align*}
\summary{Matching pair $\to$ $\delta$ = prediction $-$ label.}
\src{DLNN script WS24, Eq.~(4.70), Table~4.1.}
\end{qabox}

\begin{qabox}{Vanishing and exploding gradients. Who found it?}
\origin{asked}{Schedl committee (``backpropagation, vanishing gradients'')}
\core{The delta is multiplied by one Jacobian per layer, so over many layers it shrinks or grows exponentially.}
\explain With the sigmoid the slope is at most 0.25, so the deltas vanish. Sepp Hochreiter analysed
this in his diploma thesis in 1991, here in Linz. The fixes are ReLU or SELU, good initialization,
normalization, residual connections and, for sequences, the LSTM.
\derive
\begin{align*}
\bm{\delta}^{[1]} &= \Bigl(\prod_{l=2}^{L} \mathrm{diag}(f'(\mathbf{s}^{[l-1]}))\,W^{[l]\mathsf{T}}\Bigr)\bm{\delta}^{[L]} \why{apply the recursion $L-1$ times} \\
\lVert\bm{\delta}^{[1]}\rVert &\le \bigl(\max\lvert f'\rvert\cdot\max_l\lVert W^{[l]}\rVert\bigr)^{L-1}\lVert\bm{\delta}^{[L]}\rVert \why{norm of a product} \\
&\le \bigl(0.25\,\lVert W\rVert\bigr)^{L-1}\lVert\bm{\delta}^{[L]}\rVert \why{sigmoid: $f' \le 0.25$}
\end{align*}
If $0.25\lVert W\rVert < 1$, the delta vanishes exponentially with depth.
\summary{Product of Jacobians. Hochreiter 1991. Fix: ReLU/SELU, init, ResNet, LSTM.}
\src{DLNN script WS24, after Eq.~(4.84).}
\end{qabox}

\subsection{Gradient descent and optimizers}

\begin{qabox}{Write gradient descent. Why mini-batches?}
\origin{asked}{Schlüter, Klambauer}
\core{We move the weights a small step against the gradient of the empirical risk.}
\explain The full batch gives the exact gradient but is expensive. One sample is cheap but noisy.
A mini-batch is the compromise: a good estimate, fast on a GPU, and its noise helps to leave sharp
minima.
\derive
\begin{align*}
\mathbf{w} &\leftarrow \mathbf{w} - \eta\,\nabla_{\mathbf{w}}\Remp(\mathbf{w}) \why{full-batch gradient descent} \\
\nabla\Remp &\approx \frac{1}{\lvert B\rvert}\sum_{n\in B}\nabla L_n \why{mini-batch $B$: unbiased estimate}
\end{align*}
\summary{Step against the gradient. Mini-batch = cheap, parallel, useful noise.}
\src{DLNN script WS24, Sec.~4.10 and Ch.~7.}
\end{qabox}

\begin{qabox}{Which optimizers do you know?}
\origin{asked}{Hochreiter}
\core{Gradient descent, SGD, momentum, AdaGrad, RMSprop and Adam.}
\explain Momentum keeps a running sum of past gradients. AdaGrad and RMSprop give each weight its own
step size from the past squared gradients. Adam combines momentum and RMSprop with a bias
correction, and it is the usual default.
\derive
\begin{align*}
\mathbf{v} &\leftarrow \mu\,\mathbf{v} - \eta\,\nabla\Remp,\quad \mathbf{w} \leftarrow \mathbf{w} + \mathbf{v} \why{momentum} \\
\mathbf{m} &\leftarrow \beta_1\mathbf{m} + (1-\beta_1)\nabla,\quad \mathbf{u} \leftarrow \beta_2\mathbf{u} + (1-\beta_2)\nabla^2 \why{Adam: first and second moment} \\
\mathbf{w} &\leftarrow \mathbf{w} - \eta\,\frac{\hat{\mathbf{m}}}{\sqrt{\hat{\mathbf{u}}} + \epsilon} \why{bias-corrected step}
\end{align*}
\summary{GD, SGD, momentum, AdaGrad, RMSprop, Adam.}
\src{DLNN script WS24, Ch.~7.}
\end{qabox}

\begin{qabox}{Are there training methods other than gradient descent?}
\origin{asked}{Klambauer}
\core{Yes: closed forms, second-order methods and gradient-free methods.}
\explain Linear models have a closed form. Logistic regression can use iteratively reweighted least
squares. RProp and Quickprop use only the sign or the curvature of the gradient. Newton,
Gauss--Newton, Levenberg--Marquardt and conjugate gradients use second-order information.
Evolutionary algorithms and random search need no gradient at all.
\summary{Closed form, IRLS, RProp, Newton/LM/CG, evolutionary.}
\src{DLNN script WS24, Sec.~4.2--4.3, 7.5--7.6.}
\end{qabox}

\subsection{Activation functions, regularization, initialization}

\begin{qabox}{Which activation functions do you know? Which one is from Linz?}
\origin{asked}{Hochreiter (``which one is from Linz?'')}
\core{Sigmoid, tanh, ReLU and its variants, and ELU and SELU, which were developed at JKU.}
\explain Sigmoid and tanh are smooth but saturate. ReLU is cheap and does not saturate for positive
input, but units can die. SELU (Klambauer et al.\ 2017) makes a network self-normalizing: the mean
and variance of the activations are pulled to 0 and 1. For the output: linear for regression,
sigmoid for two classes, softmax for many.
\nums{SELU: $\lambda \approx 1.0507$, $\alpha \approx 1.6733$, fixed point $(0, 1)$.}
\summary{Linz: ELU and SELU. SELU $\to$ self-normalizing.}
\src{DLNN script WS24, Ch.~10.4 and 11.}
\end{qabox}

\begin{qabox}{Which regularization techniques do you know?}
\origin{asked}{Hochreiter}
\core{Weight decay, dropout, early stopping, data augmentation and noise.}
\explain Weight decay adds a penalty on large weights, which shrinks them in every step. Dropout
switches units off at random during training, so the network cannot rely on single units. Early
stopping ends training when the validation error starts to rise.
\derive
\begin{align*}
\tilde{R}(\mathbf{w}) &= \Remp(\mathbf{w}) + \tfrac{\lambda}{2}\lVert\mathbf{w}\rVert^2 \why{weight decay} \\
\nabla\tilde{R} &= \nabla\Remp + \lambda\,\mathbf{w} \why{gradient of the penalty} \\
\mathbf{w} &\leftarrow (1 - \eta\lambda)\,\mathbf{w} - \eta\,\nabla\Remp \why{the weights shrink each step}
\end{align*}
\summary{Weight decay, dropout, early stopping, augmentation, noise.}
\src{DLNN script WS24, Ch.~8.}
\end{qabox}

\begin{qabox}{Why does initialization matter? What is batch normalization?}
\origin{predicted}{DLNN script Ch.~9--10}
\core{Random weights break the symmetry, and their scale must keep the variance constant through the layers.}
\explain If all weights are equal, all units learn the same thing. If the scale is wrong, the
signal and the deltas vanish or explode. Batch normalization standardizes each pre-activation over
the mini-batch and then rescales it with learned parameters.
\derive
\begin{align*}
s &= \sum_{j=1}^{n_{\mathrm{in}}} w_j\,a_j \why{one pre-activation} \\
\mathrm{Var}(s) &= n_{\mathrm{in}}\,\mathrm{Var}(w)\,\mathbb{E}[a^2] \why{independent, zero-mean weights} \\
\mathrm{Var}(w) &= \frac{1}{n_{\mathrm{in}}} \why{keep $\mathrm{Var}(s) = \mathbb{E}[a^2]$} \\
\mathrm{Var}(w) &= \frac{2}{n_{\mathrm{in}}} \why{ReLU halves $\mathbb{E}[a^2]$ (He)}
\end{align*}
\summary{Break symmetry. Keep variance constant. BN = standardize, then learn scale and shift.}
\src{DLNN script WS24, Ch.~9--10.}
\end{qabox}

\subsection{Convolutional networks}

\begin{qabox}{What is a CNN and why does it work for images?}
\origin{predicted}{DLNN script Ch.~6}
\core{Small filters with shared weights slide over the image, so a pattern is found wherever it appears.}
\explain Images have local structure, and the same pattern can be anywhere. Shared weights give few
parameters and translation equivariance. Pooling makes the map smaller and adds some invariance.
Deeper layers see larger parts of the image.
\nums{Output size $\lfloor (n + 2p - k)/s \rfloor + 1$ for input $n$, kernel $k$, padding $p$,
stride $s$.}
\summary{Local filters, shared weights, pooling.}
\src{DLNN script WS24, Ch.~6.}
\end{qabox}

\begin{qabox}{ResNet, and U-Net for segmentation.}
\origin{asked}{U-Net asked by Klambauer in Life Science; ResNet predicted from the SS25 script}
\core{ResNet adds the input of a block to its output, so gradients pass through very deep networks.}
\explain The Jacobian of a residual block is the identity plus a small part, so it does not shrink
the delta. U-Net is an encoder and a decoder with skip connections. It predicts one class per pixel
and copies fine detail across the skips.
\derive
\begin{align*}
\mathbf{a}^{[l+1]} &= \mathbf{a}^{[l]} + F(\mathbf{a}^{[l]}) \why{residual block} \\
\frac{\partial\mathbf{a}^{[l+1]}}{\partial\mathbf{a}^{[l]}} &= I + \frac{\partial F}{\partial\mathbf{a}^{[l]}} \why{Jacobian stays close to $I$}
\end{align*}
\summary{ResNet: identity + residual. U-Net: encoder, decoder, skips.}
\src{DLNN I+II script SS25, ResNet and semantic segmentation.}
\end{qabox}

\subsection{Recurrent networks and LSTM}

\begin{qabox}{Explain the forward pass of a recurrent network. How is it trained?}
\origin{asked}{J.~Küng committee, Lehner (``main idea of RNNs'')}
\core{An RNN reads a sequence step by step and keeps a hidden state that carries the past.}
\explain The same weights are used at every step. It is trained with backpropagation through time:
unfold the network over the sequence and backpropagate as in a deep network with shared weights.
The number of parameters does not grow with the length of the sequence.
\derive
\begin{align*}
\mathbf{s}(t) &= W\mathbf{x}(t) + R\,\mathbf{a}(t-1),\quad \mathbf{a}(t) = f(\mathbf{s}(t)) \why{forward pass} \\
\bm\delta(t) &= \mathrm{diag}\bigl(f'(\mathbf{s}(t))\bigr)\Bigl(V^{\mathsf{T}}\frac{\partial L(t)}{\partial\hat{\mathbf{y}}(t)} + R^{\mathsf{T}}\bm\delta(t+1)\Bigr) \why{delta from the output and from the next step} \\
\frac{\partial L}{\partial R} &= \sum_t \bm\delta(t)\,\mathbf{a}(t-1)^{\mathsf{T}} \why{shared weights: sum over time}
\end{align*}
\summary{Same weights every step. Hidden state = memory. BPTT = backprop on the unfolded net.}
\src{LSTM lecture WS24 (Hochreiter).}
\end{qabox}

\begin{qabox}{Why do plain RNNs fail on long sequences, and how does LSTM fix it?}
\origin{predicted}{LSTM lecture (Hochreiter); vanishing gradients were asked}
\core{The delta is multiplied by the recurrent Jacobian at every step, so it vanishes; the LSTM cell is updated by addition instead.}
\explain The LSTM (Hochreiter and Schmidhuber 1997) has a memory cell with a constant error
carousel. Gates decide what to write, what to keep and what to read. With the forget gate near one,
the error flows back through time without shrinking.
\derive
\begin{align*}
\mathbf{c}(t) &= \mathbf{f}(t)\odot\mathbf{c}(t-1) + \mathbf{i}(t)\odot\mathbf{z}(t) \why{additive cell update} \\
\mathbf{h}(t) &= \mathbf{o}(t)\odot\tanh\bigl(\mathbf{c}(t)\bigr) \why{output gate reads the cell} \\
\frac{\partial\mathbf{c}(t)}{\partial\mathbf{c}(t-1)} &= \mathrm{diag}\bigl(\mathbf{f}(t)\bigr) \why{direct path through the cell} \\
&\approx I \why{forget gate near one: error carried unchanged}
\end{align*}
\summary{Additive cell = constant error carousel. Gates: input, forget, output.}
\src{LSTM lecture WS24, slides 6--10.}
\end{qabox}

\begin{qabox}{Bidirectional RNNs, and one label for a whole sequence.}
\origin{asked}{Schlüter (``RNN from the future'')}
\core{A bidirectional RNN runs one RNN forward and one backward and joins their states.}
\explain So every step sees the past and the future. It needs the whole sequence, so it is not for
real-time use. For one label per sequence we use the last hidden state, or we pool all hidden
states and put a classifier on top.
\summary{Forward + backward, concatenate. Sequence label: last state or pooling.}
\src{LSTM lecture WS24.}
\end{qabox}

\subsection{Attention and Transformers}

\begin{qabox}{Explain attention and the Transformer.}
\origin{asked}{Rekabsaz, Lehner, Klambauer (self-attention matrices)}
\core{Every token takes a weighted mean of all tokens, and the weights come from query times key.}
\explain Each token makes a query, a key and a value by linear maps. The softmax of query times key
gives the weights, and the output mixes the values with them. A Transformer block is multi-head
self-attention plus an MLP, each with a residual connection and layer normalization.
\derive
\begin{align*}
Q &= XW_Q,\ K = XW_K,\ V = XW_V \why{linear maps of the tokens} \\
\mathrm{Att}(Q,K,V) &= \mathrm{softmax}\Bigl(\frac{QK^{\mathsf{T}}}{\sqrt{d_k}}\Bigr)V \why{weights per row sum to one} \\
\mathrm{Var}(\mathbf{q}\cdot\mathbf{k}) &= \sum_{m=1}^{d_k}\mathrm{Var}(q_m k_m) = d_k \why{independent entries of variance 1} \\
\mathrm{Var}\bigl(\mathbf{q}\cdot\mathbf{k}/\sqrt{d_k}\bigr) &= 1 \why{so the softmax does not saturate}
\end{align*}
\summary{Query meets key $\to$ softmax weights $\to$ mix the values.}
\src{LSTM lecture WS24, slides 12--15.}
\end{qabox}

\begin{qabox}{Transformers are permutation-equivariant. Why is that a problem?}
\origin{asked}{Lehner}
\core{Self-attention sees the tokens as a set, so the order of words is lost.}
\explain If we shuffle the tokens, the outputs are shuffled the same way. So we add the position to
each token: fixed sinusoids, learned vectors, or relative encodings such as RoPE.
\summary{Attention sees a set. Add position information.}
\src{Past-student report; LSTM lecture WS24.}
\end{qabox}

\begin{qabox}{Compare RNN and Transformer in cost.}
\origin{asked}{Lehner ($O(N)$ against $O(N^2)$)}
\core{An RNN is linear in the sequence length but sequential; a Transformer is quadratic but parallel.}
\explain In a Transformer every token attends to every token, so time and memory grow with $N^2$.
But all tokens are processed at once, and long-range paths are one step long. LSTM-type models such
as xLSTM are attractive again because inference is linear in $N$.
\nums{RNN $O(N)$, Transformer $O(N^2)$.}
\summary{RNN: linear, sequential. Transformer: quadratic, parallel.}
\src{Past-student report; LSTM lecture WS24.}
\end{qabox}

\subsection{Classical machine learning in the same slot}

\begin{qabox}{Explain the support vector machine.}
\origin{asked}{Kofler, March 2026 committee}
\core{The SVM finds the separating hyperplane with the largest margin.}
\explain The margin is $2/\lVert\mathbf{w}\rVert$, so maximizing it means minimizing
$\lVert\mathbf{w}\rVert^2$. The soft margin allows some points inside the margin through slack
variables and a penalty $C$. In the dual only dot products appear, and the points with
$\alpha_n > 0$ are the support vectors.
\derive
\begin{align*}
&\min_{\mathbf{w},b,\bm\xi}\ \tfrac12\lVert\mathbf{w}\rVert^2 + C\sum_n\xi_n \why{large margin, few violations} \\
&\text{s.t.}\ \ y_n(\mathbf{w}^{\mathsf{T}}\mathbf{x}_n + b) \ge 1 - \xi_n,\ \ \xi_n \ge 0 \why{every point on the right side, up to slack} \\
&\max_{\bm\alpha}\ \sum_n\alpha_n - \tfrac12\sum_{n,m}\alpha_n\alpha_m y_n y_m\,\mathbf{x}_n^{\mathsf{T}}\mathbf{x}_m \why{dual, $0 \le \alpha_n \le C$, $\sum_n\alpha_n y_n = 0$}
\end{align*}
\summary{Max margin. Dual has only dot products.}
\src{ML: Supervised Techniques, SVM unit.}
\end{qabox}

\begin{qabox}{What is a kernel and the kernel trick?}
\origin{asked}{Kofler}
\core{A kernel is a dot product in a feature space, computed without building the features.}
\explain The SVM dual needs only dot products. We replace them by a kernel and get a non-linear
classifier at the cost of a linear one. Examples are the polynomial kernel and the Gaussian kernel.
\derive
\begin{align*}
k(\mathbf{x},\mathbf{x}') &= (\mathbf{x}^{\mathsf{T}}\mathbf{x}')^2 \why{polynomial kernel in 2D} \\
&= x_1^2x_1'^2 + 2x_1x_2x_1'x_2' + x_2^2x_2'^2 \why{multiply out} \\
&= \phi(\mathbf{x})^{\mathsf{T}}\phi(\mathbf{x}'),\quad \phi(\mathbf{x}) = (x_1^2, \sqrt2\,x_1x_2, x_2^2) \why{a dot product of features}
\end{align*}
\summary{Kernel = dot product in feature space, without computing the features.}
\src{ML: Supervised Techniques, kernels unit.}
\end{qabox}

\begin{qabox}{What is the VC dimension?}
\origin{asked}{Kofler}
\core{It is the largest number of points that a model class can label in every possible way.}
\explain It measures capacity. Linear classifiers in $d$ dimensions have VC dimension $d+1$. In
generalization bounds the gap between training and test error grows with the VC dimension and
shrinks with the number of samples.
\nums{Linear classifier in $\mathbb{R}^d$: $d+1$.}
\summary{Largest set it can shatter. Linear: $d+1$.}
\src{ML: Supervised Techniques; statistical learning theory.}
\end{qabox}

\begin{qabox}{Likelihood, maximum likelihood, and its problems.}
\origin{asked}{Widmer (several students)}
\core{Maximum likelihood chooses the parameters that make the observed data most probable.}
\explain We maximize the log-likelihood, because a sum is easier than a product. With little data
it overfits: an event never seen gets probability zero. A prior (MAP) or regularization fixes
this. Squared loss is maximum likelihood under Gaussian noise.
\derive
\begin{align*}
p(y\mid\mathbf{x}) &= \mathcal{N}\bigl(y;\,g(\mathbf{x};\mathbf{w}),\,\sigma^2\bigr) \why{Gaussian noise model} \\
-\log\prod_n p(y_n\mid\mathbf{x}_n) &= N\log\bigl(\sqrt{2\pi}\,\sigma\bigr) + \frac{1}{2\sigma^2}\sum_n\bigl(y_n - g(\mathbf{x}_n;\mathbf{w})\bigr)^2 \why{log of the Gaussian} \\
\arg\max_{\mathbf{w}}\ \text{likelihood} &= \arg\min_{\mathbf{w}}\ \sum_n\bigl(y_n - g(\mathbf{x}_n;\mathbf{w})\bigr)^2 \why{first term has no $\mathbf{w}$}
\end{align*}
\summary{Most likely parameters. Overfits with little data $\to$ prior.}
\src{DLNN script WS24, Eqs.~(4.8)--(4.11).}
\end{qabox}

\begin{qabox}{Derive the REINFORCE loss. (Asked by Prof.\ Hochreiter.)}
\origin{asked}{Hochreiter}
\core{The log-derivative trick moves the gradient inside the expectation, so we can estimate it from sampled episodes.}
\explain We want the largest expected return. The environment does not depend on the policy
parameters, so only the policy terms survive. As a loss we minimize minus log-probability times
return, with a baseline to reduce the variance.
\derive
\begin{align*}
J(\theta) &= \int p_\theta(\tau)\,G(\tau)\,\mathrm{d}\tau \why{expected return over trajectories} \\
\nabla_\theta J &= \int p_\theta(\tau)\,\nabla_\theta\log p_\theta(\tau)\,G(\tau)\,\mathrm{d}\tau \why{$\nabla p = p\,\nabla\log p$} \\
\log p_\theta(\tau) &= \log p(s_0) + \sum_t\bigl[\log\pi_\theta(a_t\mid s_t) + \log p(s_{t+1}\mid s_t,a_t)\bigr] \why{trajectory probability} \\
\nabla_\theta J &= \mathbb{E}_\tau\Bigl[\sum_t\nabla_\theta\log\pi_\theta(a_t\mid s_t)\,G_t\Bigr] \why{dynamics do not depend on $\theta$} \\
L(\theta) &= -\sum_t\log\pi_\theta(a_t\mid s_t)\,(G_t - b) \why{loss with baseline $b$}
\end{align*}
\summary{Log-derivative trick. Loss = $-\log\pi(a\mid s)\times$ return.}
\src{Deep Reinforcement Learning VL, policy gradient unit.}
\end{qabox}

\begin{qabox}{How do you evaluate a classifier?}
\origin{asked}{R.~Küng committee (AUC), Mara committee (confusion matrix)}
\core{With the confusion matrix, precision, recall and the area under the ROC curve.}
\explain Precision: of what I called positive, how much is right. Recall: of all positives, how
many I found. The ROC curve plots the true-positive rate against the false-positive rate over all
thresholds. The AUC is the probability that a random positive gets a higher score than a random
negative.
\nums{Precision $TP/(TP+FP)$. Recall $TP/(TP+FN)$. AUC 0.5 = chance, 1 = perfect.}
\summary{Precision: how clean. Recall: how complete. AUC: ranking quality.}
\src{ML: Supervised Techniques.}
\end{qabox}

\begin{qabox}{How else can you measure robustness, besides adversarial attacks?}
\origin{asked}{Klambauer}
\core{Test under distribution shift and check how sure the model is.}
\explain Use corrupted inputs, new domains or data from another time. Check calibration and
uncertainty with ensembles or MC dropout, and the sensitivity of the output to the input. My thesis
has a direct example: a recording made under one boundary condition fails under another, because
the regime changed.
\summary{Corruptions, domain shift, calibration, sensitivity, out of distribution.}
\src{DLNN script WS24, Sec.~4.8; thesis, Section 6.4.}
\end{qabox}
